% Created 2019-02-25 Mon 17:53
\documentclass[11pt]{article}
\usepackage[utf8]{inputenc}
\usepackage[T1]{fontenc}
\usepackage{fixltx2e}
\usepackage{graphicx}
\usepackage{longtable}
\usepackage{float}
\usepackage{wrapfig}
\usepackage{rotating}
\usepackage[normalem]{ulem}
\usepackage{amsmath}
\usepackage{textcomp}
\usepackage{marvosym}
\usepackage{wasysym}
\usepackage{amssymb}
\usepackage{hyperref}
\tolerance=1000
\usepackage{csquotes}
\usepackage{graphicx}
\usepackage{amsmath}
\usepackage{mathtools}
\usepackage{algpseudocode}
\usepackage[margin=0.75in]{geometry}
\usepackage{nth}
\usepackage{tikz}
\usetikzlibrary{positioning,chains,fit,shapes,shapes.misc,calc,arrows,patterns,decorations.pathreplacing}
\usepackage{ccicons}
\usepackage{calc}
\usepackage{fancyhdr}
\usepackage{fancyvrb}
\usepackage{hyperref}
\usepackage{nopageno}
\usepackage{enumitem}
\usepackage{multicol}
\usepackage{braket}
\usepackage[linguistics]{forest}
\algnewcommand\algorithmicforeach{\textbf{for each}}
\algdef{S}[FOR]{ForEach}[1]{\algorithmicforeach\ #1\ \algorithmicdo}
\setlength{\columnseprule}{1pt}
\newcommand{\bonusmarksamt}[1]{\textbf{[#1~BONUS~marks]}}
\newcommand{\bonusmarkamt}[0]{\textbf{[1~BONUS~mark]}}
\newcommand{\marksamt}[1]{\textbf{[#1~marks]}}
\newcommand{\markamt}[0]{\textbf{[1~mark]}}
\newcommand{\fillinMC}[1]{\fillinMCmath{\mbox{#1}}}
\newcommand{\fillinMCmath}[1]{\begin{tikzpicture}\draw circle [radius=0.5em];\end{tikzpicture}\ #1}
\newcommand{\fillinMCsoln}[1]{\fillinMCmathsoln{\mbox{#1}}}
\newcommand{\fillinMCmathsoln}[1]{\begin{tikzpicture}\draw[black, fill=blue] circle [radius=0.5em];\end{tikzpicture}\ #1}
\newcommand{\fillinMCAll}[1]{\fillinMCAllmath{\mbox{#1}}}
\newcommand{\fillinMCAllmath}[1]{\begin{tikzpicture}\draw (0,0) rectangle (1em,1em);\end{tikzpicture}\ #1}
\newcommand{\fillinMCAllsoln}[1]{\fillinMCAllmathsoln{\mbox{#1}}}
\newcommand{\fillinMCAllmathsoln}[1]{\begin{tikzpicture}\draw[black, fill=blue] (0,0) rectangle (1em,1em);\end{tikzpicture}\ #1}
\newcommand{\fillinblank}[1]{\fillinblankmath{\mbox{#1}}}
\newcommand{\fillinblankmath}[1]{\begingroup\setlength{\fboxsep}{1em}\setlength{\fboxrule}{2pt}\fbox{\LARGE\phantom{#1}}\endgroup}
\newcommand{\fillinblanksoln}[1]{\fillinblankmathsoln{\mbox{#1}}}
\newcommand{\fillinblankmathsoln}[1]{\begingroup\setlength{\fboxsep}{1em}\setlength{\fboxrule}{2pt}\fbox{{#1}}\endgroup}
\newcommand{\answerbox}[1][3\baselineskip]{%
\noindent\setlength{\fboxrule}{2pt}\framebox[0.7\linewidth]{%
\raisebox{2pt}[0pt][#1]{}%
}\par\medskip%
}
\usepackage{collcell,array}
\newcommand{\shiftdown}[1]{\smash{\raisebox{-.5\normalbaselineskip}{#1}}}
\newcolumntype{C}{>{\collectcell\shiftdown}c<{\endcollectcell}}
\newcommand{\shiftup}[1]{\smash{\raisebox{.5\normalbaselineskip}{#1}}}
\def\sequence#1#2{\dsequence{#1}{1}{#2}}
\def\dsequence#1#2#3{#1_{#2}, \ldots, #1_{#3}}
\newtheorem{fact}{Fact}
\def\imax{$\texttt{S}_{\max}$}
\def\iprevmax{previous\_$\texttt{S}_{\max}$}
\def\deleted{\texttt{covered}}
%%% Adding Colour to Questions and Answers
\usepackage{color}
%Questions
\def\question#1{\red{#1}}
% Answers
%Answer in short form
\def\ans#1{\par\blu{Answer: #1}}%With header of Answer
\def\soln#1{\par\blu{#1}} %Without header of Answer
% Colors
\definecolor{blu}{rgb}{0,0,0.5}
\def\blu#1{{\color{blu}#1}}
\definecolor{gre}{rgb}{0,.3,0}
\def\gre#1{{\color{gre}#1}}
\definecolor{red}{rgb}{0.5,0.0,0}
\def\red#1{{\color{red}#1}}
\def\norm#1{\|#1\|}
%%% End for Colours
\usepackage{collcell,array}
\newcolumntype{C}{>{\collectcell\shiftdown}c<{\endcollectcell}}
%\fancyhead{}\fancyfoot[C]{\tiny This work is licensed under a \href{http://creativecommons.org/licenses/by/4.0/}{Creative Commons Attribution 4.0 International License. \ccby}\\ For license purposes, the author is the University of British Columbia.}\pagestyle{fancyplain}
\makeatletter
\def\@myproof#1[#2]{\trivlist \item [\hskip\labelsep{\bf #1 {#2}: }]}
\def\@begproof#1{\@ifnextchar[{\@myproof{#1}}{\@myproof{#1}[\unskip$\,$]}}
\newenvironment{solution}{\vskip -10pt\leavevmode\@begproof{Solution}}{\endtrivlist}
\def\sopt{S_{\textit{\scriptsize opt}}}
\makeatother
\fancyhead{}\fancyfoot[C]{\tiny This work is licensed under a \href{http://creativecommons.org/licenses/by-nc/4.0/}{Creative Commons Attribution-NonCommercial 4.0 International License. \ccbync}\\ For license purposes, the author is the University of British Columbia.}\pagestyle{fancyplain}
\usepackage{mdframed}
\BeforeBeginEnvironment{minted}{\begin{mdframed}}
\AfterEndEnvironment{minted}{\end{mdframed}}
\date{}
\title{CPSC 320 2018W2: Assignment 4}
\hypersetup{
  pdfkeywords={},
  pdfsubject={},
  pdfcreator={Emacs 25.2.1 (Org mode 8.2.10)}}
\begin{document}

\maketitle
\vspace{-0.5in}

This assignment gives you practice at using recurrences to analyze
algorithm runtimes. It also gives you practice at using recurrences as
a means to design efficient algorithms, by inductively describing
quantities that we'd like to compute.

Please follow the guidelines given in Assignment 1 for submission to
Gradescope, and for group collaboration. Remember to provide short
justifications for your answers. Submit by the deadline \textbf{Monday 
March 11, 2019 at 10PM}.

Use \LaTeX to prepare your answers. Easiest will be to use the .tex
file provided.  Please enclose each paragraph of your solution in
\verb~\soln{Your solution here...}~.  \soln{Your solution will then appear in dark blue}, making it a lot easier for TAs to find the parts that you wrote.

\section{Computer scientist makes \$0.38 a day -- find out how!!!}
\label{sec-1}

A popular pastime of computer scientists (especially academics) is
fantasizing about how to use our knowledge to get rich playing the
stock market. The first advice one receives about playing with stocks
is ``buy low, sell high''.  Easier said than done.

Suppose we have an array $S=[1..n]$ of integer values where $n \ge 1$ and $S[i]$
represents the price of a stock on date $i$. Given this historic price
data, the stock market problem is to (retrospectively) find the best dates to buy and sell
the stock so as to maximize profit. If $i$ and $j$ are the buying and
selling dates (with $i \le j$), then $[i,j]$ should be a contiguous
interval with the highest jump in price, and $S[j] - S[i]$ is the
profit.  (You can buy and sell stock on the same date, with a profit of
zero.)

\vspace{.1in}

\noindent
\textbf{Example:} The $S[1..13]= [5,10,3,4,9,21,13,1,12,20,15,17,4]$ has a
maximum jump of $20-1 = 19$ if purchased on date $i=8$ and sold on date $j=10$.

\begin{enumerate}
\item \question{Describe a brute-force algorithm that identifies the best buy and sell dates, and the corresponding profit.}
\soln{
  \begin{tabbing}
  xxxxxxx\=xxx\=xxx\=xxx\=xxx\=xxxx\= \kill
  \> {\bf Algorithm} Buy-low-Sell-high($S$) \\
  \>\> $n = |S|$ \\
  \>\> $maxProfit = -\infty$ \\
  \>\> $buy = -1$ \\
  \>\> $sell = -1$ \\
  \>\> for i = 1 to n: \\
  \>\>\> for j = i to n: \\
  \>\>\>\> $profit = S[j] - S[i]$ \\
  \>\>\>\> If $profit > maxProfit:$ \\
  \>\>\>\>\> $maxProfit = profit$ \\
  \>\>\>\>\> $buy = i$ \\
  \>\>\>\>\> $sell = j$ \\
  \>\> Return $(buy, sell)$
  \end{tabbing}
}

\item \question{What is the running time of your brute force algorithm?}
\soln{
  The running time for this brute force algorithm would be $O(n^2)$. Because we have two nested for loops, 
  which grow with $n$.
}

\item A divide-and-conquer approach for the stock market problem can consider three cases:

\begin{enumerate}
\item The best buy and sell dates are both at most $\lfloor n/2 \rfloor$.
\item The best buy and sell dates are both greater than $\lfloor n/2 \rfloor$.
\item The best buy date is at most $\lfloor n/2 \rfloor$ and the best sell date
      is greater than $\lfloor n/2 \rfloor$.
\end{enumerate}

In order to determine the solution in case (c), \question{what specific pieces of information do you need to know about the left half of the array, i.e., $S[1..\lfloor n/2 \rfloor]$, and the right half, $S[\lfloor n/2 \rfloor+1.. n]$?}
\soln{
  For getting the solution on case (c) we would need to get the \emph{argmin} value of 
  $S[1..\lfloor n/2 \rfloor]$ for the best buy day and the \emph{argmax} value of the 
  $S[\lfloor n/2 \rfloor+1 .. n]$ for the best sell day.
}

\item \question{Write a recursive divide and conquer algorithm for the stock market problem.} (It may be handy for the algorithm to return
the information needed about subproblems, as documented in
your answer from part 3, in addition to the buy date, sell date and profit.)
\soln{
  \begin{tabbing}
  xxxxxxx\=xxx\=xxx\=xxx\=xxx\=xxxx\= \kill
  \> {\bf Algorithm} BlSh($S$, $start$, $end$) \\
  \>\> $n = end - start$ \\
  \>\> If $n = 0 :$ \\
  \>\>\> return $(S[start], start, start, start, start)$\\
  \>\> Else \\
  \>\>\> // Get the profit of the left and right side \\
  \>\>\> $profitL$, $buyL$, $sellL$, $minL$, $maxL$ = \texttt{BlSh}($S$, $start$, $start + \lfloor n/2 \rfloor$)\\
  \>\>\> $profitR$, $buyR$, $sellR$, $minR$, $maxR$ = \texttt{BlSh}($S$, $start + \lfloor n/2 \rfloor + 1$, $end$)\\ \\
  \>\>\> // Get the min and max value of $S[start....end]$ \\ 
  \>\>\> $amin$ = $\texttt{argmin}(S, minL, minR)$ \\
  \>\>\> $amax$ = $\texttt{argmax}(S, maxL, maxR)$ \\ \\
  \>\>\> // Calculate the profit based on the $min,max$ values \\
  \>\>\> $profit$ = $S[maxR] - S[minL]$ \\ \\
  \>\>\> // Return the maximum profit, best buy day, best sell day, minimum price, maximum price \\
  \>\>\> If $profit > profitL \text{ \&\& } profit > profitR :$ \\
  \>\>\>\> return $(profit, minL, maxR, amin, amax)$\\
  \>\>\> Else if $profitL > profit \text{ \&\& } profitL > profitR :$ \\
  \>\>\>\> return $(profitL, buyL, sellL, amin, amax)$\\
  \>\>\> Else $:$\\
  \>\>\>\> return $(profitR, buyR, sellR, amin, amax)$\\
  \end{tabbing}
}

\item \question{Give a recurrence relation for the runtime of your algorithm, and solve it to get a big-$O$ bound on the running time.}
\soln{
  \[
   \mbox{BlSh}(n) = \left \{\begin{array}{ll}
      c_1, & \mbox{if start = end } \\[.3in]
      \text{BlSh}(n/2) + \text{BlSh}(n/2) + c_2, & \mbox{otherwise} \\
      \end{array} \right.
   \]
  The recursion tree would be:
  \begin{center}
    \begin{forest}
      [$T(n)$
        [$T(n/2)$
          [$T(n/4)$]
          [$T(n/4)$]
        ]
        [$T(n/2)$
          [$T(n/4)$]
          [$T(n/4)$]
        ]
      ]
    \end{forest}
  \end{center}
  We are dividing on each level the size in half, so our algorithm will hit the base case on level 
  $i$ when $1 = \frac{n}{2^i}$, so all our leaf nodes will be on level $i = \log(n)$, thus we 
  have $\log(n)$ levels. We can also notice that on each level $i$ we are spending $2^ic_2$ time. \\ \\
  We know that we will have $n$ leafs reaching the base case with $T(1)$, and on the base case we 
  spend $c_1$ time, so for all our base cases we will have $c_1n$ time. \\ \\
  So we have to add the time at each level $i$ wich is $2^ic_2$. And we have $\log(n)-1$ levels without 
  the leaf nodes. So we have $c_2\sum_{i=0}^{\log(n) - 1}2^i$, using the geometric series formula to 
  solve that summatory, $c_2(\frac{1-2^{\log(n)}}{1-2})$ which yields $c_2n$.  \\ \\
  The running time of the algorithm would be the time of the leafs + the internal nodes. So the 
  algorithm will have $c_1n + c_2n$ running time which is $O(n)$.
}
\end{enumerate}
\section{Medians and finding the kth largest element}
\label{sec-2}
We saw in the Divide and Conquer worksheet that the QuickSelect
algorithm finds the median in $O(n)$ time, \textbf{if} we assume that we "get
lucky" with the choice of pivot at every call to QuickSelect. By getting
lucky on a call of size $n$, we mean that the rank of the pivot lies
in the range $[n/4,3n/4]$, and so the size of the next recursive call is at
most $3n/4$.  However in the worst case (when we are not lucky),
the pivot might have rank 1 or rank $n$, in which
case the algorithm takes $\Theta(n^2)$ time.

Better algorithms for choosing the pivot don't rely on luck, but
guarantee that the pivot partitions a problem of size $n$ into subproblems of
size at most $7n/10$. One such algorithm carefully chooses $\lceil n/5 \rceil$
elements in $O(n)$ time, and then sets the pivot to be the median
of these elements using a recursive call.  Exactly how the $\lceil n/5 \rceil$
elements are chosen is not important for this problem.  The rest of
the algorithm is identical to our QuickSelect algorithm. Here is
pseudocode:

\begin{tabbing}
xxxxxxx\=xxx\=xxx\=xxx\=xxx\=xxxx\= \kill
\> Algorithm  BetterQuickSelect($A[1..n], k$) \\
\>\> // Return the $k$th largest element of A  \\
\>\> // Precondition: $|A| = n \ge k$, and all elements of $A$ are unique \\ \\

\>\> // Select the pivot element $p$ \\
\>\> Create array $A'$ of $\lceil n/5 \rceil$ "carefully choosen" elements of $A$ // this takes $O(n)$ time \\
\>\> Let $p =$ BetterQuickSelect($A', \lceil$ |A'|/2 $\rceil$) \\ \\
\>\> Let Lesser  $=$ all elements from $A$ less than $p$ \\ 
\>\> Let Greater $=$ all elements from $A$ greater than $p$ \\
\>\> If |Greater| $= k - 1$ \\
\>\>\>  Return $p$ \\
\>\> Else If |Greater| $> k - 1$ \\
\>\>\>  Return BetterQuickSelect(Greater, $k$) \\
\>\> Else  // |Greater| $< k - 1$ \\
\>\>\>  Return BetterQuickSelect(Lesser, $k - |\mbox{Greater}| - 1$)
\end{tabbing}

\begin{enumerate}
\item 
Let $T'(n)$ be the running time of BetterQuickSelect.
\question{Write a recurrence for $T'(n)$.} Use
the facts that the pivot partitions a problem of size $n$ into subproblems of size at most 7n/10
and that the time to create list $A'$ is $O(n)$. \\
\soln{
  \[
   T'(n) = \left \{\begin{array}{ll}
      c, & \mbox{if |A| = 1 } \\[.3in]
      n + T'(n/5) + T'(7n/10), & \mbox{otherwise} \\
      \end{array} \right.
  \]
}

\item \question{Solve your recurrence of part 1 to show that BetterQuickSelect takes time $O(n)$ in the worst case.}
\soln{
  The recursion tree would be:
  \begin{center}
    \begin{forest}
      [$T(n)$
        [$T(n/5)$
          [$T(n/25)$]
          [$T(7n/20)$]
        ]
        [$T(7n/10)$
          [$T(7n/50)$]
          [$T(49n/100)$]
        ]
      ]
    \end{forest}
  \end{center}
  By analizing the recursion tree we can notice that at the $2$nd level we are spending $cn\frac{9}{10}$ time 
  because each call is divided into two calls of sizes $\frac{n}{5}$, $\frac{7n}{10}$ This is a pattern 
  that repeats on following $i$ level with a running time of $cn\frac{9}{10}^i$. So the total cost, 
  would be represented by the following sum $cn\sum_{i=0}^{\infty} \frac{9}{10}^i$, that can be solved 
  by the geometric series, $cn\sum_{i=0}^{\infty} \frac{9}{10}^i = cn\frac{1}{1-\frac{9}{10}} = 10cn$, 
  thus the total running time of BetterQuickSelect is $O(n)$
}
\end{enumerate}

\section{More Expressions of 1's (Adapted from Jeff Erickson)}
\label{sec-3}

As in the tutorial of week 8, we can represent positive integers in
terms of 1's, $+$, $\times$, and parentheses "(" and ")".  For
example, all of the following expressions represent the integer 14:
\[
\begin{array}{l}
1 + 1 + 1 + 1 + 1 + 1 + 1 + 1 + 1 + 1 + 1 + 1 + 1 + 1 + 1 + 1 \\
((1 + 1) \times (1 + 1 + 1 + 1 + 1)) + ((1 + 1) \times (1 + 1)) \\
(1 + 1) \times (1 + 1 + 1 + 1 + 1 + 1 + 1) \\
(1 + 1) \times (((1 + 1 + 1) \times (1 + 1)) + 1) 
\end{array}
\]
Given a positive integer $n$, we'd like to find the minimum number of
1's in such an expression.  Let $M(n)$ denote this minimum.  For
example, the fourth expression above is such an expression for 14, and
so $M(14) = 8$.

\begin{enumerate}
\item \question{Use the recurrence for $M(n)$ from the solutions to the tutorial of week 8, plus memoization, to develop an efficient recursive algorithm to compute $M(n)$.}
\soln{
  \begin{tabbing}
  xxxxxxx\=xxx\=xxx\=xxx\=xxx\=xxxx\= \kill
  \> {\bf Algorithm} $M(n, Soln)$ \\
  \>\> // Soln is the array with solutions, at the begining it is filled with $n$ $-1$s \\
  \>\> $Min = \infty$ \\
  \>\> If $n = 1 :$ \\
  \>\>\> return $1$ \\
  \>\> If $Soln[n] != -1 :$ \\
  \>\>\> return $Soln[n]$ \\
  \>\> For $i = 1$ to $n-1$ $:$ \\
  \>\>\> $Min = $ \texttt{min}($Min$, $M(n-i)+M(i)$) \\
  \>\>\> If $(i+1) < n/2 \text{ \&\& } (i+1)|n : $ \\
  \>\>\>\> $Min = $ \texttt{min}($Min$, $M(n/(i+1))+M(i+1)$) \\
  \>\> $Soln[n] = Min$ \\
  \>\> return $Min$ \\
  \end{tabbing}
}

\item \question{What is the running time of your memoized algorithm?}
\end{enumerate}

\section{Max-Weight Independent Sets}
\label{sec-4}
Let $G=(V,E,w())$ be an undirected graph, with a weight $w(v)$
associated with each node $v\in V$.  An \emph{independent set} of $G$ is a
subset of nodes, none of which are connected by an edge of the
graph. A \emph{max-weight} independent set is an
independent set such that the sum of the weights of the nodes in the
set is as large as possible. 
Let MIS($G$) be the weight of a max-weight independent set for graph $G$.
Finding max-weight independent sets and their corresponding weights
can be done efficiently for simple types of graphs, although no efficient
algorithm is known for general graphs. Here, you'll consider graphs
that are paths.

For $0 \le i \le n$, let $P = ([1..n],E,w)$ be a weighted path, that is $E = \{(1,2),
(2,3), \ldots, (i,i+1), \ldots (n-1,n)\}$.  
For example, in the following path, the weights are the numbers drawn
inside the nodes.  

\vspace{.1in}
\hspace{.75in}
\includegraphics[width=5in]{../figs/path-graph.png}

\vspace{.1in}

\noindent
Let MIS$^+(P,i)$ be the max weight of an independent
set of $P$ involving only nodes from $[1..i]$ and including $i$.  
Similarly, let MIS$^-(P,i)$ be the max weight of an independent set
involving only nodes from $[1..i-1]$ (node $i$ is excluded from the independent
set). We can write MIS$^+(P,i)$ as:
\[
\mbox{MIS}^+(P,i) = \left \{\begin{array}{ll}
                   0, & \mbox{if } i = 0 \\
                   w(i) +  \mbox{MIS}^-(P,(i-1)), & \mbox{if } i \ge 1.
                   \end{array} \right.
\]

\begin{enumerate}
\item \question{Write a recurrence for $\mbox{MIS}^-(P,i)$, $i\ge 0$:}

\[
   \mbox{MIS}^-(P,i) = \left \{\begin{array}{ll}
                   \rule{.3in}{0.4pt}, & \mbox{if } \rule{.5in}{0.4pt} \\[.3in]
                   \max \{  \rule{1in}{0.4pt}, \rule{1in}{0.4pt} \}, & \mbox{if } \rule{.5in}{0.4pt} \\
                   \end{array} \right.
   \]

\item 
We can express MIS$(P)$ as $\max \{\mbox{MIS}^+(P,n), \mbox{MIS}^-(P,n) \}$. 
\question{Use this to write an algorithm that computes MIS($P$) in $O(n)$ time.}  
Your algorithm should use memoization and recursion to compute both
MIS$^+(P,n)$ and MIS$^-(P,n)$, and use these quantities to compute
MIS($P$). (No justification is needed, just the algorithm.)

\item \question{Using a recursion tree, explain why the running time of your algorithm of part 3 is $\Theta(n)$.}
\end{enumerate}

\section{More Max-Weight Independent Sets}
\label{sec-5}
This problem concerns independent sets on "rail line" graphs. These
undirected graphs have $2n$ nodes, labeled $[1,i]$, $[2,i]$, $1 \le i
\le n$. A function $w()$ assigns a weight to each node. There are
edges between $[1,i]$ and $[2,i]$ for each $i$, as well as edges
between $[1,i]$ and $[1,i+1]$ for $1 \le i \le n-1$ and between
$[2,i]$ and $[2,i+1]$ for $1 \le i \le n-1$, for a total of $2(n-1) +
n = 3n-2$ edges.

\question{Develop recurrences that could be used to efficiently compute MIS($G$) when the input $G$ is a rail line graph.} You
can express MIS($G$) in terms of other quantities of your
choosing that you can define, and then provide recurrences for these quantities.
You do \emph{not} need to describe an algorithm to compute MIS($G$).
% Emacs 25.2.1 (Org mode 8.2.10)
\end{document}
